\bigskip
\begin{center}\textls[120]{c) Die Theorie des agM. bis zum 30. Mai 1799.}\end{center}
\indent Zwischen den auf die lemniskatischen Funktionen und andere Gegenstände
bezüglichen Aufzeichnungen finden sich im \textit{Leiste} bei den Seiten 25, 26, 37,\par
\clearpage\SourceLeaf{127}\ArithmeticHead{42}{L}{L. SCHLESINGER, ÜBER GAUSS’ FUNKTIONENTHEORETISCHE ARBEITEN.}
\noindent 48, 49, 56, 83, 88 Notizen über das agM., die Werke X, 1, S. 177—180 abgedruckt
sind. Diese Notizen sind irreguläre und stammen allem Anschein
nach aus der zweiten Hälfte des Jahres 1799, werden aber vielleicht zum Teil
als Niederschlag aus älteren Aufzeichnungen anzusehen sein, die nicht mehr
erhalten sind. — Der älteste auf das agM. bezügliche Zettel (Ff) ist Werke
X, 1, S. 172, 173 abgedruckt und trägt in lateinischer Sprache die Überschrift
»Über das, wenn man so sagen darf, arithmetisch-geometrische Mittel«. Er
enthält zunächst die Reihe
\[
\begin{aligned}
\operatorname{Tm}(1+x)=1+z={}&1+\frac12x-\frac1{16}xx+\frac1{32}x^3-\frac{21}{1024}x^4+\frac{31}{2048}x^5\\
&-\frac{195}{16384}x^6+\frac{305}{32768}x^7+\ldots,
\end{aligned}
\]
wo \(\operatorname{Tm}(1+x)\) das agM. (Terminus medius) zwischen \(1\) und \(1+x\) bedeutet,
und die Umkehrung derselben:
\[
1+x=1+2z+\frac12zz-\frac14z^3+\frac3{16}z^4-\frac5{32}z^5+\frac{23}{256}z^6-\frac5{128}z^7+\ldots.
\]
Es wird dann (vergl. die Bemerkungen Werke X, 1, S. 257—260) eine eindeutige
Funktion \(\varepsilon^{*u*}\) eingeführt, die zu einer Uniformisierung der Beziehung
\(1+z=\operatorname{Tm}(1+x)\) in der Form
\[
1+x=\varepsilon^{*u*},\quad1+z=\varepsilon^{*\frac{u}{2}*}
\]
dienen sollte. Über die Mängel, mit denen dieser Ansatz behaftet ist, vergleiche
man die erwähnten Bemerkungen; hier genügt es festzustellen, dass
\textsc{Gauss} eine Uniformisierung der durch das agM. gesetzten Beziehung schon
zu einer Zeit ins Auge gefasst hatte, als er noch mit den primitiven Methoden
des formalen Operierens mit Potenzreihen arbeitete, bei denen er Inversion,
Einsetzen von einer Reihe in eine andere u. dgl. durch das Berechnen einiger
erster Glieder zu erledigen suchte. Hält man sich vor Augen, dass diese
Methoden in der Theorie der lemniskatischen Funktion schon zu Anfang des
Jahres 1797 durch die numerisch induktive, mehr quantitative Methode im
Sinne \textsc{Eulers} abgelöst worden sind, so wird man veranlasst, den Ursprung
dieses Zettels auf eine sehr frühe Zeit zu setzen. Eine obere Zeitgrenze ergibt
sich daraus, dass auf der letzten Seite der Scheda Aa mit tändelnder
Schrift »Medium Arith. G.« und ebenda auch »Hermannstadt« steht; nun hat
\textsc{Gauss} am 31. Oktober 1798 einen Brief von \textsc{Wolfgang Bolyai} empfangen\par
\clearpage\SourceLeaf{128}\ArithmeticHead{43}{R}{III c. DAS AGM. BIS ZUM MAI 1799.}
\HistoricalFoot{}{\(6^*\)}
\noindent (\textit{Briefwechsel}, S. 8, 9), in dem dieser von seiner siebenbürgischen Heimat
schreibt, und damit stimmt die Tatsache, dass \textsc{Gauss} die nächste Scheda (Ab)
im November 1798 begonnen hat. Er hat sich also Ende Oktober 1798 mit
dem agM. beschäftigt, und unser Zettel ist sicher vor dieser Zeit geschrieben.

Die aus dem Anfang des Jahres 1799 stammenden, auf das agM. bezüglichen
Aufzeichnungen in der Scheda Ab (Werke X, 1, S. 174—176) stehen
allerdings methodisch auch nicht viel höher, sie enthalten aber immerhin
schon einiges Numerische, worunter besonders die auf das agM. zwischen 1
und \(\sqrt2\) Bezug habenden Zahlen hervorzuheben wären; so steht z. B. Werke
X, 1, S. 174 unten die Zahl \(0{,}19814\), die gleich \(M(1,\sqrt2)-1\) ist. — Was in
den erwähnten Leistenotizen aus ältern Aufzeichnungen entnommen, und was
erst in der zweiten Hälfte des Jahres 1799 gefunden worden ist, lässt sich
nur vermuten, jedenfalls zeigt der Inhalt dieser Notizen, dass sie erst nach
der Tagebuchaufzeichnung 98. vom 30. Mai 1799 geschrieben sein können.
Diese Aufzeichnung lautet wie folgt: »Dass das agM. zwischen 1 und \(\sqrt2\) gleich
\(\frac{\pi}{\varpi}\) ist, haben wir bis zur elften Dezimalziffer bestätigt; wenn dies bewiesen
sein wird, so ist damit sicher ein wahrhaft neues Feld der Analysis erschlossen«.

Eine Aufzeichnung, wo diese numerische Vergleichung angegeben ist, hat
sich nicht vorgefunden. In bezug auf die Berechnung der Grösse \(\frac{\pi}{\varpi}\) haben
wir gesehen, dass \textsc{Gauss} nach den verschiedenartigsten Methoden ihren Wert
mit immer wachsender Genauigkeit bestimmt hat; zuletzt hat er noch mit
Hilfe der Reihen (4) die Grösse \(\sqrt{\frac{\varpi}{\pi}}\) auf 26 Dezimalstellen berechnet und,
indem er diese Zahl ins Quadrat erhebt, \(\frac{\varpi}{\pi}\) selbst bestimmt (siehe genaueres
über diese Rechnung Werke X, 1, S. 538). In bezug auf die Berechnung von
\(M(1,\sqrt2)\) haben wir aus der damaligen Zeit nur die oben erwähnte Aufzeichnung
in der Scheda Ab, aber in der nachgelassenen Abhandlung über das
agM. von 1800 (Werke III, S 364) wird dieser Wert auf 19 Dezimalstellen
angegeben, und es ist mit Sicherheit anzunehmen, dass \textsc{Gauss} diese sehr einfache
Rechnung schon 1799 ausgeführt haben wird.

Welche Bedeutung \textsc{Gauss} der Bemerkung vom 30. Mai beigelegt hat,
geht unter anderem daraus hervor, dass er dieselbe an \textsc{Johann Friedrich Pfaff}
in Helmstedt brieflich mitgeteilt hat. Der Brief selbst ist — wie alle Briefe
von \textsc{Gauss} an \textsc{Pfaff}, mit Ausnahme eines einzigen (Werke X, 1, S. 250) —\par
\clearpage\SourceLeaf{129}\ArithmeticHead{44}{L}{L. SCHLESINGER, ÜBER GAUSS’ FUNKTIONENTHEORETISCHE ARBEITEN.}
\noindent nicht aufzufinden, aber aus der im Gaussarchiv vorhandenen Antwort von
\textsc{Pfaff} vom 24. November 1799 (Werke X, 1, S. 232, 233) geht mit Sicherheit
hervor, dass \textsc{Gauss}, als er \textsc{Pfaff} diese Mitteilung machte, noch nicht im Besitz
eines analytischen Beweises war. Ja man kann aus der Bemerkung, die
\textsc{Pfaff} in jenem Briefe macht (a. a. O., S. 232): »so leicht das arithmetisch-harmonische
Mittel zu finden und zu beweisen ist\footnote{Für dieses ist, wenn \(a,b\) die Ausgangszahlen sind \(a_1=\frac12(a+b),\ b_1=2ab/(a+b)\) usw., also
\(a_1b_1=ab\) und folglich, wenn man mit \(A\) den Grenzwert bezeichnet, \(A^2=ab\).\hfill\textsc{Schlesinger}.}, so wenig scheint bei
der ersten Ansicht klar zu sein, wie das agM. zu bestimmen sei« schliessen,
dass \textsc{Gauss} damals überhaupt noch keine Beziehung zwischen dem agM. und
irgendwelcher »klassischen Grösse«, also auch nicht dem Ellipsenquadranten,
gekannt hat, da er sonst \textsc{Pfaff} sicher davon Mitteilung gemacht hätte, wenn
auch vielleicht in der Form eines »pli cacheté«, wie in dem folgenden Briefe
in bezug auf das elementare Mittel, das zu dem Grenzwert \(\sqrt{a^2-b^2}/\operatorname{arc\,cos}\frac ba\)
führt (siehe \textsc{Pfaffs} Antwortschreiben, Werke X, 1, S. 234). — Hiernach wird
man also annehmen dürfen, dass \textsc{Gauss}, als er durch die Beobachtung vom
30. Mai 1799 veranlasst wurde, sich erneut mit dem agM. zu beschäftigen,
über den elementaren Algorithmus und das, was in den Aufzeichnungen
Werke X, 1, S. 172—173 enthalten ist, hinausgehend nur noch die Beziehung
zu den Potenzreihen, deren Exponenten eine arithmetische Progression zweiter
Ordnung bilden, gekannt haben dürfte, und dass er erst im Sommer oder
Spätherbst 1799 zu weiteren Erkenntnissen fortgeschritten ist.
